\documentclass[10pt,a4paper]{amsart}
\usepackage[margin=19mm]{geometry}
\usepackage{amsmath,amssymb,mathtools}
\usepackage[scr=rsfs,cal=euler]{mathalfa}
\usepackage{xfrac}
\usepackage[unicode,hidelinks,bookmarks=false]{hyperref}
\newcommand{\ie}{i.e.\ }
\newcommand{\cf}{cf.\ }
\newcommand{\resp}{respectively\ }
\newcommand{\Subsection}{\subsection}
\providecommand{\nobreakdash}{\nobreak\hbox{-}}
\setlength{\emergencystretch}{2em}
\multlinegap0pt
\usepackage[shortlabels]{enumitem}
\usepackage{amssymb}
\newtheorem{prop}{Proposition}
\newtheorem{defi}{Definition}
\newtheorem{thm}{Theorem}
\newcommand{\ctB}[6]{{\mathsf B}_{#1, #2, #3}#4#5#6}
\newcommand{\ctC}[6]{{\mathsf C}_{#1, #2, #3}#4#5#6}
\newcommand{\ctD}[6]{{\mathsf D}^{#1}_{#2, #3, #4}#5#6}
\newcommand{\ctI}[2]{{\mathsf I}_{#1}#2}
\newcommand{\ctK}[4]{{\mathsf K}_{#1, #2}#3#4}
\newcommand{\ctS}[6]{{\mathsf S}_{#1, #2, #3}#4#5#6}
\newcommand{\ctW}[4]{{\mathsf W}_{#1, #2}#3#4}
\newcommand{\lbd}[3] {
\lambda{#1}\hspace{-.75mm}:\hspace{-.75mm}{#2}\hspace{.35mm}.\hspace{.35mm}{#3}
}
\newcommand{\mltup}{\boldsymbol\lambda_{\param}}
\newcommand{\param}{\theta}
\newcommand{\reduces}[2]{\rightarrow_{#1}^{#2}}
\newcommand{\transreduces}[2]{\twoheadrightarrow_{#1}^{#2}}
\newcommand{\typd}[2] {
#1\hspace{-.75mm}:\hspace{-.75mm}#2}
\title{Simply-typed constant-domain modal lambda calculus~I: distanced beta reduction\texorpdfstring{\\}{} and combinatory logic}
\author{Sean Walsh}
\date{Source: arXiv:2410.17463v6 (CC BY 4.0)}
\begin{document}
\maketitle

\noindent\emph{Source and licence.} Simply-typed constant-domain modal lambda calculus~I: distanced beta reduction\texorpdfstring{\\}{} and combinatory logic. Version arXiv:2410.17463v6 (CC BY 4.0), distributed by arXiv under the Creative Commons Attribution 4.0 International licence, http://creativecommons.org/licenses/by/4.0/, as recorded in the arXiv licence metadata for this exact version and shown on the abstract page https://arxiv.org/abs/2410.17463v6. Full licence notice: \texttt{licenses/arxiv-2410.17463-CC-BY-4.0.txt}.\par\smallskip

\noindent\textit{Editorial scope.} Counted unit: the article's thm prop:combinatoryeffect (source0.tex lines 1026--1032). The article's complete original proof of it is delivered with it (source0.tex lines 1033--1103, one \texttt{proof} environment of 71 lines). No mathematics is written, repaired or re-derived: the statement and the proof are the article's own lines, in the article's own notation, and no retained line is substituted or re-flowed.  Retained uncounted as the source-local context the proof needs: lines 319--321; lines 446--466; lines 963--1007.  Nothing was omitted from any retained span, so the package carries no omission record.  No retained line contains a citation, so the wrapper carries no bibliography and no bibliography entry.  No retained line contains a figure environment, an image-inclusion command or drawing code, so no image is delivered and the package declares no asset dependency.  Editorial formatting only, in addition: an AMS \texttt{amsart} preamble that restates the article's own declarations and macros used by the retained lines, namely the environments \texttt{prop}, \texttt{defi}, \texttt{thm} on separate counters, together with the standard packages those lines need.  The retained environments print in this document as Proposition 1, Definition 1, Theorem 1 in order of appearance, whereas the article prints the same items with the numbering of its own full document.  The complete native source of the article as distributed in the arXiv e-print for this exact version is bundled unchanged as \texttt{sources/167279/source0.tex}.
\begin{prop}[Terms of state type]\label{prop:state-free}
The only terms of $\mltup$ of state type are the variables and constants.
\end{prop}

\begin{defi}[Definition of $\beta$- and $\eta$-reduction in $\mltup$]\label{defn:betareduction}
We say $\big( \lbd{\vec{x}}{\vec{B}}{\lbd{v}{A}{L}}\big) \vec{M} N \reduces{\beta}{\param} \big( \lbd{\vec{x}}{\vec{B}}{L[v:=N]} \big) \vec{M}$ if each of the following conditions holds
\begin{enumerate}[leftmargin=*]
    \item\label{defn:betareduction:1} $\typd{N}{A}$ is free for $\typd{v}{A}$ in $\typd{L}{C}$;
    \item\label{defn:betareduction:3} the variables in $\typd{\vec{x}}{\vec{B}}$ are not free in $\typd{N}{A}$;
    \item \label{defn:betareduction:2} the variables in $\typd{\vec{x}}{\vec{B}}, \typd{v}{A}$ are pairwise distinct.
\end{enumerate}

The \emph{distance} of the $\beta$-reduction is the length of vector $\typd{\vec{x}}{\vec{B}}$.

We use $\beta_0$ for distance zero $\beta$-reduction.

We say that an instance of $\beta$-reduction is \emph{regular} if
\begin{enumerate}[leftmargin=*]\setcounter{enumi}{3}
    \item\label{defn:betareduction:5} the only variable, if any, of state type in the tuple $\typd{\vec{x}}{\vec{B}}\equiv \typd{x_1}{B_1}, \ldots, \typd{x_n}{B_n}$ is the first one $\typd{x_1}{B_1}$.
\end{enumerate}

We use $\beta_r$ for regular $\beta$-reduction.

Finally, we define $\lbd{x}{A}{Mx} \reduces{\eta}{\param} M$ when $\typd{x}{A}$ not free in $\typd{M}{A\rightarrow B}$.
\end{defi}

\begin{defi}[Typed combinator terms of $\mltup$]\label{defn:typedcombo}
An \emph{Identity Bird term} $\ctI{A}{}$ \emph{of} $\mltup$ is a term of 
of the following form and type:
\begin{equation*}
\lbd{x}{A}{x} \; : \; A\rightarrow A
\end{equation*}
It is required that $A$ has regular type. Intuitively $\ctI{A}{}$ is the identity function on type $A$.

A \emph{Kestrel term} $\ctK{A}{B}{}{}$ \emph{of} $\mltup$  is a term of the following form and type:
\begin{equation*}
\lbd{x}{A}{\lbd{y}{B}{x}} \; : \; A  \rightarrow B\rightarrow   A
\end{equation*}
It is required that $A$ has regular type, and that $\typd{x}{A}, \typd{y}{B}$ are distinct variables. Intuitively given a value in $A$, Kestrel $\ctK{A}{B}{}{}$ returns the constant function from $B$ to $A$ with that value.

A \emph{Cardinal term} $\ctC{A}{B}{C}{}{}{}$ \emph{of} $\mltup$ is a term of the following form and type:
\begin{equation*}
\lbd{x}{A \rightarrow B\rightarrow C}{\lbd{y}{B}{\lbd{z}{A}{xzy}}} \;:\; (A\rightarrow B\rightarrow   C) \rightarrow B\rightarrow  A\rightarrow  C
\end{equation*}
It is required that $C$ has regular type, and that $\typd{y}{B}, \typd{z}{A}$ are distinct variables. Intuitively Cardinal $\ctC{A}{B}{C}{}{}{}$ takes a function $x$ of two arguments and returns the function of two arguments which permutes the two inputs. I.e. it maps function $x$ to the function $(y,z)\mapsto x(z,y)$.

A \emph{Dardinal term} $\ctD{c}{A}{B}{C}{}{}$ \emph{of} $\mltup$  is a term of the following form and type:
\begin{equation*}
\lbd{x}{A \rightarrow B\rightarrow C}{\lbd{z}{A}{xzc}} \; : \; (A\rightarrow B\rightarrow   C) \rightarrow  A\rightarrow  C
\end{equation*}
It is required that $C$ has regular type and that $B$ has state type and that $\typd{c}{B}$ is a constant.  \emph{Dardinal} is short for \emph{decorated Cardinal}. The number of Dardinals varies with the signature, and intuitively $\ctD{c}{A}{B}{C}{}{}$ takes a function $x$ of two arguments and returns the function of one argument which slots this value into the first spot and $c$ into the second spot. I.e. it maps function $x$ to the function $z\mapsto x(z,c)$.

A \emph{Starling term} $\ctS{A}{B}{C}{}{}{}$ \emph{of} $\mltup$  is a term of the following form and type:
\begin{equation*}
\lbd{x}{C\rightarrow  A\rightarrow  B}{\lbd{y}{C\rightarrow   A}{\lbd{z}{C}{xz(yz)}}} \;:\; (C\rightarrow   A\rightarrow  B) \rightarrow (C\rightarrow A)\rightarrow C\rightarrow B
\end{equation*}
It is required that $A,B$ are regular types. Like Cardinal, Starling $\ctS{A}{B}{C}{}{}{}$ permutes some of the order of the inputs, but it also is a basic example of a combinator which duplicates an input. As for its intended behaviour, as an argument of $x,y$, it is just an ``indexed''  version of functional application $z\mapsto x_z(y_z)$.

A \emph{Warbler term} $\ctW{A}{B}{}{} $ \emph{of} $\mltup$  is a term of the following form and type:
\begin{equation*}
\lbd{x}{A \rightarrow A \rightarrow B}{\lbd{y}{A}{xyy}} \; : \; (A\rightarrow  A\rightarrow  B) \rightarrow A\rightarrow  B
\end{equation*}
It is required that $B$ is of regular type. Intuitively, Warbler $\ctW{A}{B}{}{} $ takes a curried function $x$ defined on $A\times A$ and returns the diagonal function on $A$ given by $y\mapsto x(y,y)$.

A \emph{Bluebird term} $\ctB{A}{B}{C}{}{}{}$ \emph{of} $\mltup$  is a term of the following form and type: 
\begin{equation*}
\lbd{x}{B\rightarrow C}{\lbd{y}{A\rightarrow B}{\lbd{z}{A}{x(yz)}}} : (B\rightarrow C) \rightarrow (A \rightarrow B) \rightarrow A\rightarrow  C
\end{equation*}
It is required that $B,C$ are regular types; and if $A$ is regular then it is required that $\typd{x}{B\rightarrow C}, \typd{z}{A}$ are distinct variables and that $\typd{x}{B\rightarrow C}, \typd{y}{A\rightarrow B}$ are distinct variables. The Bluebird $\ctB{A}{B}{C}{}{}{}$ returns the composition of two functions $x,y$ whose domains and codomains match appropriately. 
 
\end{defi}

\begin{thm}[Combinatory behaviour in $\mltup$]\label{prop:combinatoryeffect}
Suppose that the combinators on the below left are terms of $\mltup$. For each item, suppose that the terms $P,Q,R$ are terms of $\mltup$ of the appropriate type to make the below applications well-formed. Then one has the $\beta$-reductions to the terms on the below right, and indeed all of these reductions are regular:
\begin{align*}
& \ctI{A}{P}\transreduces{\beta}{\param} P  \hspace{7mm} \ctK{A}{B}{P}{Q}\transreduces{\beta}{\param} P & \hspace{1mm}\ctC{A}{B}{C}{P}{Q}{R}\transreduces{\beta}{\param} PRQ & \hspace{5mm}\ctD{c}{A}{B}{C}{P}{R}\transreduces{\beta}{\param} PRc\\
& \hspace{3mm}\ctS{A}{B}{C}{P}{Q}{R}\transreduces{\beta}{\param} PR(QR) & \hspace{5mm} \ctW{A}{B}{P}{Q}\transreduces{\beta}{\param} PQQ  &\hspace{5mm}  \ctB{A}{B}{C}{P}{Q}{R}\transreduces{\beta}{\param} P(QR) 
\end{align*}
\end{thm}

\begin{proof}
For  Identity Bird $\ctI{A}{}$, we have $(\lbd{x}{A}{x})P\reduces{\beta}{\param} P$ since the term $\typd{x}{A}$ has no lambda abstracts.

Now we turn to the remaining cases. In this part of the proof:
\begin{itemize}[leftmargin=*]
\item We restrict attention to the case where at least one of the constituent types $A,B,C$ (or just $A,B$ in the case of Kestrel and Warbler) is a state type. For, if all of the constituent types are regular, then by inspection of the definitions, all of the bound variables in the combinator terms are regular and we can change them by $\alpha$-conversion to avoid variable capture. 
\item We repeatedly use Proposition~\ref{prop:state-free}, which says that terms of state type in $\mltup$ are either constants or variables of that very state type. 
\item We often ensure that the ``free for'' condition of $\beta$-reduction (Definition~\ref{defn:betareduction}(\ref{defn:betareduction:1})) is met by doing $\alpha$-conversion on the bound variables of the combinatory terms to ensure that they are disjoint from the free variables of the inputs. We refer to this simply as ``disjointness.''
\item As we proceed, we note that the reductions meet the regularity condition Definition~\ref{defn:betareduction}(\ref{defn:betareduction:5}), but do not say more than this since it follows clearly from the displayed instances and the case assumptions. Further, so as not to clutter the proof, we just mark regularity in the text and write the simpler $\beta$ instead of $\beta_r$. Finally, since all $\beta$-reductions of distance $\leq 1$ are trivially regular, we only need to explicitly take note of regularity for instances of distance $\geq 2$ (as a matter of fact, all $\beta$-reductions in this proof have distance $\leq 2$).
\end{itemize}

For Kestrel $\ctK{A}{B}{}{}$, suppose $A$ is a regular type, $B$ is a state type, and $\typd{P}{A}$ and $\typd{Q}{B}$ are terms of $\mltup$. We can use $\alpha$-conversion on the bound variable of type $A$ in Kestrel so that it does not appear free in $Q$. Further, $\typd{Q}{B}$ is free for Kestrel's second variable $\typd{y}{B}$ in the term $\typd{x}{A}$, since the latter has no lambda abstracts. These two points get us the following, where the first is a $\beta$-reduction of distance~$1$:
\begin{equation}
\big(\lbd{x}{A}{\lbd{y}{B}{x}}\big) P Q \reduces{\beta}{\param} (\lbd{x}{A}{x[y:=Q]})P \equiv (\lbd{x}{A}{x})P \reduces{\beta}{\param} P
\end{equation}
The second $\beta$-reduction follows just as in Identity Bird.

For Cardinal $\ctC{A}{B}{C}{}{}{}$, suppose that $C$ is a regular type and suppose we have the terms $\typd{P}{A\rightarrow B\rightarrow C}$ and $\typd{Q}{B}$ and $\typd{R}{A}$ of $\mltup$. Suppose the corresponding three bound variables of Cardinal $\ctC{A}{B}{C}{}{}{}$  are $\typd{x}{A\rightarrow B\rightarrow C}$ and $\typd{y}{B}$ and $\typd{z}{A}$. We respectively refer to these in the following discussion as the \emph{first}, \emph{second}, and \emph{third bound variables} of Cardinal (and we adopt similar conventions for the subsequent combinator terms). 

There are three cases to consider. 

If $A$ is a regular type and $B$ is a state type, then since the first and third bound variables are of regular type, we may change them by $\alpha$-conversion so that they do not appear free in $P$ or $R$; and they do not appear free in $\typd{Q}{B}$ since $B$ is of state type. Then $\typd{Q}{B}$ is free for $\typd{y}{B}$ in $\lbd{z}{A}{xzy}$. Further, the first bound variable does not appear free in $\typd{Q}{B}$ since $B$ is of state type. This gives us the first step in the following, which is a $\beta$-reduction of distance~$1$:
\begin{align}
  (\lbd{x}{A\rightarrow B\rightarrow C}{\lbd{y}{B}{\lbd{z}{A}{xzy}}})PQR  &    \reduces{\beta}{\param} \;  (\lbd{x}{A\rightarrow B\rightarrow C}{\lbd{z}{A}{xzQ}})PR \notag \\
  \reduces{\beta}{\param} \;   (\lbd{z}{A}{PzQ})R  & \reduces{\beta}{\param} \; PRQ 
\end{align}
The second $\beta$-reduction follows by disjointness: the third bound variable does not appear free in $P$ by the previous $\alpha$-conversion; and $\typd{Q}{B}$ does not contain any bound variables since $B$ is of state type . The third $\beta$-reduction follows since the displayed free occurrence of $\typd{z}{A}$ is the only free occurrence in $PzQ$, since by previous $\alpha$-conversion it does not appear free in $P$, and since $\typd{Q}{B}$ is of state type $B$.

If $A$ is a state type and $B$ is a regular type, then since the first and the second bound variables are of regular type, we may change them by $\alpha$-conversion so that they do not appear free in $P$ or $Q$; and they do not occur free in $\typd{R}{A}$ since $A$ is of state type. Since $xzy$ contains no lambda abstracts, one has that $\typd{R}{A}$ is free for $\typd{z}{A}$ in $xzy$. This gives us the first step in the following, which is a regular $\beta$-reduction of distance~$2$:
\begin{align}
 (\lbd{x}{A\rightarrow B\rightarrow C}{\lbd{y}{B}{\lbd{z}{A}{xzy}}}) PQR & \reduces{\beta}{\param}\; (\lbd{x}{A\rightarrow B\rightarrow C}{\lbd{y}{B}{xRy}}) PQ \notag \\
\reduces{\beta}{\param}\; (\lbd{y}{B}{PRy}) Q & \reduces{\beta}{\param}\; PRQ
\end{align}
The second $\beta$-reduction follows by disjointness: the second bound variable does not appear free in $P$ by the previous $\alpha$-conversion; and $\typd{R}{A}$ does not contain any bound variables since $A$ is of state type. The third $\beta$-reduction follows since the displayed free occurrence of $\typd{y}{B}$ is the only free occurrence in $PRy$, since by previous $\alpha$-conversion it does not appear free in $P$, and since $\typd{R}{A}$ is of state type.

If $A,B$ are both state types, then since the first bound variable is of regular type it does not appear free in $\typd{Q}{B}$ or $\typd{R}{A}$ since these are of state type. Since the last two bound variables $\typd{y}{B}, \typd{z}{A}$ of Cardinal are distinct by definition (cf. Definition~\ref{defn:typedcombo}) by $\alpha$-conversion we can assume that if $\typd{Q}{B}$ is a variable then it is the second bound variable $\typd{y}{B}$. This implies that $\typd{Q}{B}$ is free for $\typd{y}{B}$ in $\lbd{z}{A}{xzy}$. Then we have the following, where the first is a $\beta$-reduction of distance 1:
\begin{align}
(\lbd{x}{A\rightarrow B\rightarrow C}{\lbd{y}{B}{\lbd{z}{A}{xzy}}}) PQR  & \reduces{\beta}{\param}\; (\lbd{x}{A\rightarrow B\rightarrow C}{\lbd{z}{A}{xzQ}}) PR \notag \\
 \reduces{\beta}{\param}\; (\lbd{x}{A\rightarrow B\rightarrow C}{xRQ}) P & \reduces{\beta}{\param}\; PRQ
\end{align}
The second $\beta$-reduction is of distance~1 and follows because of the fact that $\typd{R}{A}$ is free for $\typd{z}{A}$ in $xzQ$ since the term $xzQ$ has no lambda abstracts in it since $\typd{Q}{A}$ is of state type. The third $\beta$-reduction follows since the displayed free occurrence of $\typd{x}{A\rightarrow B\rightarrow C}$ is the only free occurrence in $xRQ$, due to $\typd{Q}{B}$ and $\typd{R}{A}$ being of state type.


For Dardinal $\ctD{c}{A}{B}{C}{}{}$, suppose that $C$ is a regular type and $A,B$ are state types with $\typd{c}{B}$ a constant. Suppose we have the terms $\typd{P}{A\rightarrow B\rightarrow C}$ and $\typd{R}{A}$. Then $\typd{R}{A}$ is free for $\typd{z}{A}$ in $xzc$ since this term has no lambda abstracts. And the first bound variable $\typd{x}{A\rightarrow B\rightarrow C}$ does not appear free in $\typd{R}{A}$ since $A$ is of state type. Then we have the following, where the first $\beta$-reduction is of distance~1:
\begin{equation}
\big( \lbd{x}{A\rightarrow B\rightarrow C}{\lbd{z}{A}{xzc}}\big) P R \reduces{\beta}{\param} (\lbd{x}{A\rightarrow B\rightarrow C}{ x R c})P 
 \reduces{\beta}{\param} PRc
\end{equation}
The second $\beta$-reduction happens since the displayed free occurrence of $\typd{x}{A\rightarrow B\rightarrow C}$ in $xRc$ is the only occurrence since $\typd{R}{A}$ is of state type.

For Starling $\ctS{A}{B}{C}{}{}{}$, suppose that $A,B$ are regular types and $C$ is a state type. Starling's first two bound variables are of regular type, and by $\alpha$-conversion we may assume that they do not appear free in $P,Q$; and they do not appear free in $\typd{R}{C}$ since it is of state type. Since the term $xz(yz)$ contains no lambda abstracts, one has that $\typd{R}{C}$ is free for $\typd{z}{C}$ in $xz(yz)$. Then we have the following, where the first application of $\beta$ is regular of distance~2:
\begin{align}
& (\lbd{x}{C\rightarrow A\rightarrow B}{\lbd{y}{C\rightarrow A}{\lbd{z}{C}{xz(yz)}}})PQR \notag \\ & \reduces{\beta}{\param} \; (\lbd{x}{C\rightarrow A\rightarrow B}{\lbd{y}{C\rightarrow A}{xR(yR)}})PQ \notag \\
& \reduces{\beta}{\param} \; (\lbd{y}{C\rightarrow A}{PR(yR)})Q \reduces{\beta}{\param} \; PR(QR) 
\end{align}
The second $\beta$-reduction follows by disjointness: the second bound variable does not appear free in $P$ by the previous $\alpha$-conversion; and $\typd{R}{C}$ does not contain any bound variables since $C$ is of state type. The third $\beta$-reduction follows since the displayed free occurrence of $\typd{y}{C\rightarrow A}$ is the only free occurrence in $PR(yR)$, since by previous $\alpha$-conversion it does not appear free in $P$, and it does not appear free in $\typd{R}{C}$ since this is of state type $C$.


For Warbler $\ctW{A}{B}{}{}$, suppose $B$ is a regular type, and $A$ is a state type, and $\typd{P}{A\rightarrow A\rightarrow B}$ and $\typd{Q}{A}$ are terms. Since Warbler's first bound variable is of regular type, by $\alpha$-conversion we may assume that it does not appear free in $P$; and it does not appear free in $\typd{Q}{A}$ since this is of state type. Since the term $xyy$ does not contain any lambda abstracts, one has that $\typd{Q}{A}$ is free for $\typd{y}{A}$ in $xyy$. Then we have the following, where the first instance of $\beta$-reduction is of distance~$1$:
\begin{equation}
(\lbd{x}{A\rightarrow A\rightarrow B}{\lbd{y}{A}{xyy}})PQ  \reduces{\beta}{\param} \; (\lbd{x}{A\rightarrow A\rightarrow B}{xQQ})P\reduces{\beta}{\param} \;  PQQ
\end{equation}
The last $\beta$-reduction follows since the displayed instance of $\typd{x}{A\rightarrow A\rightarrow B}$ is the only free instance in $xQQ$ since $\typd{Q}{A}$ is of state type. 

For Bluebird $\ctB{A}{B}{C}{}{}{}$, suppose that $B,C$ are regular types, and $A$ is a state type. Since Bluebird's first two bound variables are of regular type, by $\alpha$-conversion we may assume that they do not appear free in $P,Q$; and they do not appear free in $\typd{R}{A}$ since it is of state type. Further, $\typd{R}{A}$ is free for $\typd{z}{A}$ in $x(yz)$ since the term $x(yz)$ contains no lambda abstracts. Then we have the following, where the first instance of $\beta$ is regular of distance~2:
\begin{align}
 (\lbd{x}{B\rightarrow C}{\lbd{y}{A\rightarrow B}{\lbd{z}{A}{x(yz)}}})PQR & \reduces{\beta}{\param} \; (\lbd{x}{B\rightarrow C}{\lbd{y}{A\rightarrow B}{x(yR)}})PQ \notag \\
\reduces{\beta}{\param} \; (\lbd{y}{A\rightarrow B}{P(yR)})Q & \reduces{\beta}{\param} \;  P(QR)
\end{align}
The second $\beta$-reduction follows by disjointness: the second bound variable does not appear free in $P$ by the previous $\alpha$-conversion; and $\typd{R}{A}$ does not contain any bound variables since $A$ is of state type. The third $\beta$-reduction follows since the displayed free occurrence of $\typd{y}{B}$ is the only free occurrence in $P(yR)$, since by previous $\alpha$-conversion it does not appear free in $P$, and it does not appear free in $\typd{R}{A}$ since this is of state type $A$.
\end{proof}

\end{document}
